\section{Cartesian coordinates}

To describe points on a line, in a plane, or in space, we use coordinates. The most common coordinate system is the Cartesian coordinate system.

\begin{center}
\begin{tabular}{@{}lll@{}}
\toprule
Setting & Coordinates & Example \\
\midrule
Line & \(x\) & \(2\) \\
Plane & \((x,y)\) & \((2,3)\) \\
Space & \((x,y,z)\) & \((1,-2,5)\) \\
\bottomrule
\end{tabular}
\end{center}

This horizontal summary is often more useful than writing the three cases separately. It shows at once how the number of coordinates grows with the dimension.

\begin{figure}[htbp]
\centering
\input{figures/mathematical_foundations/cartesian_point}
\caption{A point represented by Cartesian coordinates in the plane.}
\label{fig:mathematical-foundations-cartesian-point}
\end{figure}

Coordinates are naturally connected with the way we describe location. A coordinate system gives a mathematical way of saying where something is.

\begin{remarkbox}
Coordinates are not the point itself. Coordinates depend on a choice of origin, axes, orientation, and units. The same point can have different coordinates in different coordinate systems.
\end{remarkbox}

Geometry concerns objects such as points, lines, planes, and curves. Coordinates are a method of describing these objects using numbers.
